\documentclass[a4paper,12pt]{article}

\usepackage[portuguese]{babel}
\usepackage[utf8]{inputenc}
\usepackage[T1]{fontenc}
\usepackage{amsmath,amssymb,amsthm}
\usepackage{geometry}
\usepackage{listings}
\usepackage{xcolor}
\usepackage{hyperref}

\geometry{margin=2.5cm}

\title{Construções Categóricas em Categorias Algébricas\\
\large com exemplos em MATLAB/Octave}
\author{}
\date{}

\definecolor{codegray}{rgb}{0.95,0.95,0.95}
\definecolor{darkgreen}{rgb}{0,0.45,0}
\definecolor{darkblue}{rgb}{0,0,0.6}

\lstset{
    backgroundcolor=\color{codegray},
    basicstyle=\ttfamily\small,
    breaklines=true,
    frame=single,
    columns=fullflexible,
    keepspaces=true,
    keywordstyle=\color{darkblue},
    commentstyle=\color{darkgreen},
    showstringspaces=false
}

\begin{document}

\maketitle

\section*{Introdução}

Neste documento estudamos, para várias categorias algébricas, as construções de
produtos, coproductos, pullbacks, pushouts, equalizadores e coequalizadores.
Além disso, indicamos funtores e transformações naturais entre alguns pares de
categorias.

Como o enunciado pede também uma descrição simplificada em \texttt{MATLAB/Octave},
incluímos exemplos computacionais elementares. Esses exemplos são
\textbf{didáticos} e servem para ilustrar as construções; não implementam toda a
generalidade algébrica de forma completamente formal.

\section*{Categorias consideradas}

\begin{itemize}
    \item $\mathbf{Mag}$: magmas $(M,\ast)$.
    \item $\mathbf{DIMag}$: magmas com duas operações $(M,\ast,\circ)$.
    \item $\mathbf{Mon}$: monóides $(M,\ast,e)$.
    \item $\mathbf{OrdMon}$: monóides ordenados $(M,\ast,e,\le)$.
    \item $\mathbf{Lat}$: reticulados $(L,\wedge,\vee)$.
    \item $\mathbf{DLat}$: reticulados distributivos.
    \item $\mathbf{Bool}$: álgebras booleanas.
\end{itemize}

\section*{Representação em MATLAB/Octave}

Em \texttt{MATLAB/Octave}, representamos as estruturas por \texttt{structs}. Por
exemplo, um magma pode ser representado por um conjunto finito de elementos e uma
operação binária:

\begin{lstlisting}[language=Matlab]
% Exemplo de magma
A.elements = [0 1 2];
A.op = @(x,y) mod(x + y, 3);
\end{lstlisting}

Um monóide acrescenta um elemento neutro:

\begin{lstlisting}[language=Matlab]
M.elements = [0 1 2];
M.op = @(x,y) mod(x + y, 3);
M.e = 0;
\end{lstlisting}

Um reticulado pode ser representado por duas operações:

\begin{lstlisting}[language=Matlab]
L.elements = [0 1];
L.meet = @(x,y) min(x,y);
L.join = @(x,y) max(x,y);
\end{lstlisting}

\section*{1. Categoria \(\mathbf{Mag}\)}

\subsection*{Produto}

Sejam $(A,\ast_A)$ e $(B,\ast_B)$ magmas. O produto é o magma
\[
A\times B
\]
com operação definida por
\[
(a,b)\ast(a',b')=(a\ast_A a',\, b\ast_B b').
\]

\paragraph{Código MATLAB/Octave:}
\begin{lstlisting}[language=Matlab]
function P = magma_product(A, B)
    P.elements = [];
    
    for i = 1:length(A.elements)
        for j = 1:length(B.elements)
            P.elements = [P.elements; A.elements(i), B.elements(j)];
        end
    end
    
    P.op = @(x,y) [A.op(x(1), y(1)), B.op(x(2), y(2))];
end
\end{lstlisting}

\subsection*{Coproducto}

O coproducto existe e é o \emph{magma livre} gerado por $A\sqcup B$, impondo as
relações já válidas em $A$ e em $B$. Denota-se informalmente por
\[
A\ast_{\mathbf{Mag}} B.
\]

\paragraph{Comentário computacional:}
uma implementação completa exigiria representar \emph{termos formais} e impor
relações de equivalência; isso é possível, mas já não é simples em
\texttt{MATLAB/Octave}. Por isso, aqui apenas indicamos a construção teórica.

\subsection*{Pullback}

Dados morfismos $f:A\to C$ e $g:B\to C$, o pullback é
\[
P=\{(a,b)\in A\times B \mid f(a)=g(b)\},
\]
com a operação herdada do produto.

\paragraph{Código MATLAB/Octave:}
\begin{lstlisting}[language=Matlab]
function P = pullback(A, B, f, g)
    P.elements = [];
    
    for i = 1:length(A.elements)
        for j = 1:length(B.elements)
            if f(A.elements(i)) == g(B.elements(j))
                P.elements = [P.elements; A.elements(i), B.elements(j)];
            end
        end
    end
    
    P.op = @(x,y) [A.op(x(1), y(1)), B.op(x(2), y(2))];
end
\end{lstlisting}

\subsection*{Pushout}

Dados $f:C\to A$ e $g:C\to B$, o pushout é obtido como quociente do coproducto
$A\ast_{\mathbf{Mag}} B$ pela menor congruência que identifica
\[
i_A(f(c)) \sim i_B(g(c)), \qquad c\in C.
\]

\paragraph{Comentário computacional:}
como no coproducto, uma implementação completa requer manipulação de classes de
equivalência e congruências.

\subsection*{Equalizador}

Dados $f,g:A\to B$, o equalizador é
\[
E=\{a\in A \mid f(a)=g(a)\}.
\]

\paragraph{Código MATLAB/Octave:}
\begin{lstlisting}[language=Matlab]
function E = equalizer(A, f, g)
    E.elements = [];
    
    for i = 1:length(A.elements)
        a = A.elements(i);
        if f(a) == g(a)
            E.elements = [E.elements a];
        end
    end
    
    E.op = A.op;
end
\end{lstlisting}

\subsection*{Coequalizador}

Dados $f,g:A\to B$, o coequalizador é o quociente
\[
B/{\sim},
\]
onde $\sim$ é a menor congruência tal que
\[
f(a)\sim g(a),\qquad a\in A.
\]

\paragraph{Código simplificado MATLAB/Octave:}
\begin{lstlisting}[language=Matlab]
function classes = coequalizer(A, f, g)
    classes = {};
    
    for i = 1:length(A.elements)
        classes{end+1} = [f(A.elements(i)), g(A.elements(i))];
    end
end
\end{lstlisting}

\paragraph{Observação:}
este código é apenas ilustrativo; não constrói o quociente completo, apenas
regista as identificações que deveriam gerar a congruência.

\section*{2. Categoria \(\mathbf{DIMag}\)}

Um objecto é $(M,\ast,\circ)$, com duas operações binárias.

\subsection*{Produto}

O produto de $(A,\ast_A,\circ_A)$ e $(B,\ast_B,\circ_B)$ é
\[
A\times B
\]
com
\[
(a,b)\ast(a',b')=(a\ast_A a',\, b\ast_B b'),
\]
\[
(a,b)\circ(a',b')=(a\circ_A a',\, b\circ_B b').
\]

\paragraph{Código MATLAB/Octave:}
\begin{lstlisting}[language=Matlab]
function P = dimag_product(A, B)
    P.elements = [];
    
    for i = 1:length(A.elements)
        for j = 1:length(B.elements)
            P.elements = [P.elements; A.elements(i), B.elements(j)];
        end
    end
    
    P.op1 = @(x,y) [A.op1(x(1), y(1)), B.op1(x(2), y(2))];
    P.op2 = @(x,y) [A.op2(x(1), y(1)), B.op2(x(2), y(2))];
end
\end{lstlisting}

\subsection*{Coproducto, Pushout e Coequalizador}

Existem e são construídos como objectos livres/quocientes por congruências,
agora compatíveis com \emph{duas} operações.

\subsection*{Pullback}

\[
P=\{(a,b)\in A\times B\mid f(a)=g(b)\},
\]
com ambas as operações herdadas.

\subsection*{Equalizador}

\[
E=\{a\in A\mid f(a)=g(a)\},
\]
subdouble magma de $A$.

\section*{3. Categoria \(\mathbf{Mon}\)}

\subsection*{Produto}

Se $(M,\ast,e_M)$ e $(N,\cdot,e_N)$ são monóides, então o produto é
\[
M\times N
\]
com
\[
(m,n)\bullet(m',n')=(m\ast m',\, n\cdot n'),
\qquad
e_{M\times N}=(e_M,e_N).
\]

\paragraph{Código MATLAB/Octave:}
\begin{lstlisting}[language=Matlab]
function P = monoid_product(M, N)
    P.elements = [];
    
    for i = 1:length(M.elements)
        for j = 1:length(N.elements)
            P.elements = [P.elements; M.elements(i), N.elements(j)];
        end
    end
    
    P.op = @(x,y) [M.op(x(1), y(1)), N.op(x(2), y(2))];
    P.e = [M.e, N.e];
end
\end{lstlisting}

\subsection*{Coproducto}

O coproducto é o produto livre de monóides,
\[
M\ast_{\mathbf{Mon}} N.
\]

\paragraph{Comentário computacional:}
em exemplos finitos simples, pode-se representar palavras alternadas; mas, em
geral, a implementação é já substancialmente mais pesada.

\subsection*{Pullback}

Dados $f:M\to P$ e $g:N\to P$,
\[
M\times_P N=\{(m,n)\in M\times N\mid f(m)=g(n)\},
\]
com operação e unidade herdadas.

\subsection*{Pushout}

É o quociente do coproducto pela menor congruência de monóides que identifica
as imagens de cada elemento do objecto intermédio.

\subsection*{Equalizador}

\[
E=\{m\in M\mid f(m)=g(m)\},
\]
submonóide de $M$.

\subsection*{Coequalizador}

Quociente de $N$ pela menor congruência de monóides que força
$f(m)\sim g(m)$ para todo $m\in M$.

\section*{4. Categoria \(\mathbf{OrdMon}\)}

Um objecto é $(M,\ast,e,\le)$, onde $(M,\ast,e)$ é monóide e a multiplicação é
monótona em cada variável.

\subsection*{Produto}

O produto de dois monóides ordenados é
\[
M\times N
\]
com ordem produto:
\[
(m,n)\le(m',n') \iff m\le_M m' \text{ e } n\le_N n',
\]
e multiplicação componente a componente.

\paragraph{Código MATLAB/Octave:}
\begin{lstlisting}[language=Matlab]
function P = ordmon_product(M, N)
    P.elements = [];
    
    for i = 1:length(M.elements)
        for j = 1:length(N.elements)
            P.elements = [P.elements; M.elements(i), N.elements(j)];
        end
    end
    
    P.op = @(x,y) [M.op(x(1), y(1)), N.op(x(2), y(2))];
    P.e = [M.e, N.e];
    P.leq = @(x,y) M.leq(x(1),y(1)) && N.leq(x(2),y(2));
end
\end{lstlisting}

\subsection*{Coproducto}

Existe, mas a sua descrição explícita é mais técnica: trata-se do monóide
ordenado livre gerado por $M\sqcup N$, impondo as relações de operação e de
ordem já existentes.

\subsection*{Pullback, Equalizador}

Definem-se como subestruturas do produto, respeitando tanto a operação como a
ordem.

\subsection*{Pushout, Coequalizador}

Obtêm-se por quocientes ordenados adequados, forçando as identificações
necessárias.

\section*{5. Categoria \(\mathbf{Lat}\)}

\subsection*{Produto}

Se $L$ e $M$ são reticulados, então
\[
L\times M
\]
com
\[
(\ell,m)\wedge(\ell',m')=(\ell\wedge\ell',\, m\wedge m'),
\]
\[
(\ell,m)\vee(\ell',m')=(\ell\vee\ell',\, m\vee m').
\]

\paragraph{Código MATLAB/Octave:}
\begin{lstlisting}[language=Matlab]
function P = lattice_product(L, M)
    P.elements = [];
    
    for i = 1:length(L.elements)
        for j = 1:length(M.elements)
            P.elements = [P.elements; L.elements(i), M.elements(j)];
        end
    end
    
    P.meet = @(x,y) [L.meet(x(1), y(1)), M.meet(x(2), y(2))];
    P.join = @(x,y) [L.join(x(1), y(1)), M.join(x(2), y(2))];
end
\end{lstlisting}

\subsection*{Coproducto}

Existe e é o reticulado livre sobre $L\sqcup M$, impondo as relações já válidas
em $L$ e em $M$.

\subsection*{Pullback}

Dados $f:L\to K$ e $g:M\to K$,
\[
P=\{(\ell,m)\in L\times M \mid f(\ell)=g(m)\},
\]
subreticulado de $L\times M$.

\subsection*{Pushout}

É o quociente do coproducto pela menor congruência que identifica as imagens
correspondentes.

\subsection*{Equalizador}

\[
E=\{x\in L\mid f(x)=g(x)\},
\]
subreticulado de $L$.

\subsection*{Coequalizador}

Quociente pela menor congruência de reticulados que identifica $f(x)$ e $g(x)$.

\section*{6. Categoria \(\mathbf{DLat}\)}

As construções são análogas às de $\mathbf{Lat}$, mas dentro da subcategoria dos
reticulados distributivos.

\subsection*{Produto}

O produto cartesiano de dois reticulados distributivos continua a ser
distributivo, logo o produto é dado por operações componente a componente.

\subsection*{Coproducto}

Existe em $\mathbf{DLat}$ e é o reticulado distributivo livre sobre
$L\sqcup M$ com as relações internas de $L$ e de $M$.

\subsection*{Pullback, Equalizador}

Continuam a ser subreticulados do produto.

\subsection*{Pushout, Coequalizador}

Continuam a ser quocientes por congruências apropriadas.

\section*{7. Categoria \(\mathbf{Bool}\)}

\subsection*{Produto}

Se $B_1$ e $B_2$ são álgebras booleanas, então
\[
B_1\times B_2
\]
com operações definidas componente a componente:
\[
(x_1,x_2)\wedge(y_1,y_2)=(x_1\wedge y_1,\, x_2\wedge y_2),
\]
\[
(x_1,x_2)\vee(y_1,y_2)=(x_1\vee y_1,\, x_2\vee y_2),
\]
\[
\neg(x_1,x_2)=(\neg x_1,\neg x_2).
\]

\paragraph{Código MATLAB/Octave:}
\begin{lstlisting}[language=Matlab]
function P = bool_product(B1, B2)
    P.elements = [];
    
    for i = 1:length(B1.elements)
        for j = 1:length(B2.elements)
            P.elements = [P.elements; B1.elements(i), B2.elements(j)];
        end
    end
    
    P.meet = @(x,y) [B1.meet(x(1), y(1)), B2.meet(x(2), y(2))];
    P.join = @(x,y) [B1.join(x(1), y(1)), B2.join(x(2), y(2))];
    P.compl = @(x) [B1.compl(x(1)), B2.compl(x(2))];
end
\end{lstlisting}

\subsection*{Coproducto, Pushout, Coequalizador}

Existem e são obtidos por construções livres e quocientes booleanos.

\subsection*{Pullback, Equalizador}

São subálgebras booleanas do produto.

\section*{8. Funtores e transformações naturais}

\subsection*{(p) \(\mathbf{DIMag}\) e \(\mathbf{Mag}\)}

Existem dois funtores covariantes naturais:
\[
U_1:\mathbf{DIMag}\to \mathbf{Mag},
\qquad
U_1(M,\ast,\circ)=(M,\ast),
\]
\[
U_2:\mathbf{DIMag}\to \mathbf{Mag},
\qquad
U_2(M,\ast,\circ)=(M,\circ).
\]

Também existe o funtor diagonal:
\[
\Delta:\mathbf{Mag}\to \mathbf{DIMag},
\qquad
\Delta(M,\ast)=(M,\ast,\ast).
\]

Temos isomorfismos naturais
\[
U_1\circ \Delta \cong \mathrm{Id}_{\mathbf{Mag}},
\qquad
U_2\circ \Delta \cong \mathrm{Id}_{\mathbf{Mag}}.
\]

\paragraph{Código MATLAB/Octave:}
\begin{lstlisting}[language=Matlab]
function M = forget_first(D)
    M.elements = D.elements;
    M.op = D.op1;
end

function M = forget_second(D)
    M.elements = D.elements;
    M.op = D.op2;
end

function D = diagonal_magma(M)
    D.elements = M.elements;
    D.op1 = M.op;
    D.op2 = M.op;
end
\end{lstlisting}

\subsection*{(q) \(\mathbf{Lat}\) e \(\mathbf{DLat}\)}

Existe o funtor de inclusão:
\[
I:\mathbf{DLat}\hookrightarrow \mathbf{Lat}.
\]

Também existe um reflector distributivo:
\[
D:\mathbf{Lat}\to \mathbf{DLat},
\]
que envia um reticulado $L$ no seu maior quociente distributivo
\[
D(L)=L/\theta_{\mathrm{dist}},
\]
onde $\theta_{\mathrm{dist}}$ é a menor congruência que força a distributividade.

Há uma transformação natural
\[
q:\mathrm{Id}_{\mathbf{Lat}} \Rightarrow I\circ D,
\]
em que $q_L$ é a projecção canónica no quociente distributivo.

\paragraph{Comentário:}
o exemplo ``$\mathbf{Lat}\to \mathbf{DLat}$ restringindo aos distributivos''
não define um funtor em toda a categoria $\mathbf{Lat}$, porque nem todo
reticulado é distributivo. O que existe naturalmente é a inclusão ou então o
quociente reflector.

\subsection*{(r) \(\mathbf{DLat}\) e \(\mathbf{Bool}\)}

Existe o funtor de esquecimento
\[
J:\mathbf{Bool}\to \mathbf{DLat},
\]
que vê uma álgebra booleana apenas como reticulado distributivo.

Se estivermos a trabalhar com reticulados distributivos limitados, podemos
definir o functor
\[
Z:\mathbf{DLat}\to \mathbf{Bool},
\]
por
\[
Z(L)=\{a\in L \mid a \text{ é complementado em }L\}.
\]

Em reticulados distributivos, o conjunto dos elementos complementados forma uma
álgebra booleana.

Há então uma transformação natural
\[
\iota: J\circ Z \Rightarrow \mathrm{Id}_{\mathbf{DLat}},
\]
cujas componentes são as inclusões do centro booleano em $L$.

\subsection*{(s) \(\mathbf{OrdMon}\) e \(\mathbf{Mon}\)}

Existe o funtor de esquecimento
\[
U:\mathbf{OrdMon}\to \mathbf{Mon},
\qquad
U(M,\ast,e,\le)=(M,\ast,e).
\]

Existe também o funtor de ordem discreta
\[
\mathrm{Disc}:\mathbf{Mon}\to \mathbf{OrdMon},
\]
definido por
\[
x\le_{\mathrm{disc}} y \iff x=y.
\]

Temos um isomorfismo natural
\[
U\circ \mathrm{Disc}\cong \mathrm{Id}_{\mathbf{Mon}}.
\]

\paragraph{Código MATLAB/Octave:}
\begin{lstlisting}[language=Matlab]
function M = forget_order(OrdM)
    M.elements = OrdM.elements;
    M.op = OrdM.op;
    M.e = OrdM.e;
end

function OrdM = discrete_order(M)
    OrdM = M;
    OrdM.leq = @(x,y) (x == y);
end
\end{lstlisting}

\section*{9. Exemplo completo em MATLAB/Octave}

Segue-se um exemplo simples de uso para o produto de magmas.

\begin{lstlisting}[language=Matlab]
A.elements = [0 1];
A.op = @(x,y) mod(x+y,2);

B.elements = [0 1];
B.op = @(x,y) mod(x*y,2);

P = magma_product(A,B);

disp(P.op([1 1],[1 0]))
\end{lstlisting}

Neste caso:
\begin{itemize}
    \item em $A$, a operação é a soma módulo $2$;
    \item em $B$, a operação é o produto módulo $2$;
    \item no produto, a operação é feita componente a componente.
\end{itemize}

\section*{10. Observações finais}

\begin{itemize}
    \item Em $\mathbf{Mag}$, $\mathbf{DIMag}$, $\mathbf{Mon}$, $\mathbf{Lat}$,
    $\mathbf{DLat}$ e $\mathbf{Bool}$, os \textbf{limites} (produtos, pullbacks,
    equalizadores) são tipicamente dados por subestruturas de produtos.

    \item Os \textbf{colimites} (coproductos, pushouts, coequalizadores) são
    construídos por objectos livres e quocientes por congruências.

    \item Em $\mathbf{OrdMon}$ o mesmo princípio continua válido, embora os
    colimites sejam tecnicamente mais delicados devido à presença simultânea da
    operação e da ordem.

    \item Os exemplos em \texttt{MATLAB/Octave} apresentados aqui são
    simplificados e servem sobretudo para ilustrar os casos finitos mais básicos.

    \item Em contexto teórico, os funtores mais naturais entre as categorias
    pedidas são:
    \[
    U_1,U_2:\mathbf{DIMag}\to \mathbf{Mag},
    \qquad
    \Delta:\mathbf{Mag}\to \mathbf{DIMag},
    \]
    \[
    I:\mathbf{DLat}\hookrightarrow \mathbf{Lat},
    \qquad
    D:\mathbf{Lat}\to \mathbf{DLat},
    \]
    \[
    J:\mathbf{Bool}\to \mathbf{DLat},
    \qquad
    Z:\mathbf{DLat}\to \mathbf{Bool},
    \]
    \[
    U:\mathbf{OrdMon}\to \mathbf{Mon},
    \qquad
    \mathrm{Disc}:\mathbf{Mon}\to \mathbf{OrdMon}.
    \]
\end{itemize}

\end{document}